\documentclass[12pt]{exam}
\usepackage{amsthm}
\usepackage{bm}
\usepackage{libertine}
\usepackage[utf8]{inputenc}
\usepackage[margin=0.5in]{geometry}
\usepackage{amsmath,amssymb}
\usepackage{multicol}
\usepackage[shortlabels]{enumitem}
\usepackage{siunitx}
\usepackage{physics}
\usepackage{booktabs}
\usepackage{graphicx}
\usepackage{pgfplots}
\usepackage{listings}
\usepackage{tikz}
\usepackage{tikz-3dplot}

\bmdefine{\ii}{i}
\bmdefine{\jj}{j}
\bmdefine{\kk}{k}
\newcommand{\te}{\theta}
\newcommand{\cC}{\mathcal{C}}
\newcommand{\word}[1]{\quad\text{#1}\quad}

\pgfplotsset{width=10cm,compat=1.9}
\usepgfplotslibrary{external}
\tikzexternalize

\newcommand{\class}{Math 261-001}
\newcommand{\examnum}{Quiz 2 Solutions}
\newcommand{\examdate}{September 4}

\begin{document}
\pagestyle{plain}
\thispagestyle{empty}

\noindent
\textbf{\class}\\
\textbf{\examnum}, \textbf{\examdate}

\vspace{10pt}
\textbf{Grading:} Students answer exactly two problems. Each selected problem
is worth 10 points, for a total of 20 points.

\small
\begin{enumerate}

\item \textbf{Expected final answers:}
\[
  \text{(a) Yes.}\qquad \text{(b) No.}
\]
\textbf{Verification:} The direction vectors satisfy
\[
  (2,4,-2)=2(1,2,-1),
\]
so the lines are parallel. They are distinct because
$(0,1,4)-(1,0,2)=(-1,1,2)$ is not a scalar multiple of $(1,2,-1)$.
Thus they do not intersect.

\textbf{Equivalent forms:} ``Parallel'' for part (a) and ``do not intersect''
or ``disjoint'' for part (b) receive full credit.

\textbf{Grading (10 points):} 4 points for recognizing the proportional
direction vectors, 3 points for showing that the lines are distinct, and
3 points for the two correct conclusions.

\item \textbf{Expected final answer:}
\[
  \vec{r}(t)=\left(5\cos t,\frac{5}{2}\sin t\right),\qquad
  \vec{v}(\pi/2)=\vec{r}'(\pi/2)=(-5,0).
\]
\textbf{Verification:} Substitution gives
\[
  (5\cos t)^2+4\left(\frac{5}{2}\sin t\right)^2
  =25(\cos^2t+\sin^2t)=25,
\]
and
\[
  \vec{v}(t)=\vec{r}'(t)
  =\left(-5\sin t,\frac{5}{2}\cos t\right),\qquad
  \vec{v}(\pi/2)=(-5,0).
\]
The parametrization starts at $(5,0)$ and moves counterclockwise.

\textbf{Equivalent forms:} Algebraically identical notation is valid.\par
Any counterclockwise parametrization is also valid when paired with its
derivative at $t=\pi/2$.

\textbf{Grading (10 points):} 6 points for the parametrization and 4 points
for the derivative and its value at $t=\pi/2$.

\item \textbf{Expected final answer:}
\[
  \vec{r}(t)=(0,t,-3t),\qquad t\in\mathbb{R}.
\]
\textbf{Verification:} The plane normals are
\[
  \vec{n}_1=(1,3,1),\qquad \vec{n}_2=(1,6,2),\qquad
  \vec{n}_1\times\vec{n}_2=(0,-1,3).
\]
Setting $y=0$ in both plane equations gives the point $P=(0,0,0)$.
Thus $\vec{r}(t)=t(0,-1,3)$ parametrizes the intersection; replacing $t$ by
$-t$ gives the displayed answer.

\textbf{Equivalent forms:} Any affine parametrization of the same line
$x=0$, $z=-3y$ receives full credit. In particular, any nonzero scalar
multiple of the direction vector $(0,1,-3)$ is valid.

\textbf{Grading (10 points):} 4 points for a correct cross product, 2 points
for a point on both planes, and 4 points for a complete parametric equation.

\item \textbf{Expected final answer:}
\[
  t=2,\qquad \vec{r}(2)=(2,4,4).
\]
\textbf{Verification:} At the film, $t+t^2+2t=10$, so
$t^2+3t-10=(t+5)(t-2)=0$.
Because $t\geq0$, the only allowed value is $t=2$, and
$\vec{r}(2)=(2,4,4)$. The point satisfies $2+4+4=10$.

\textbf{Equivalent forms:} The time must equal $2$; the point may be written
as an ordered triple or with labeled coordinates.

\textbf{Grading (10 points):} 4 points for the plane-intersection equation,
3 points for the allowed parameter value, and 3 points for the point of
intersection.

\end{enumerate}

\end{document}
